% ═══════════════════════════════════════════════════════════════════════
% OPERA — Analysis Plan
% Outcome-guided Patient Embedding for Registry Adaptation
%
% Insert into Overleaf as a standalone document or \input{} into main.
% Requires: booktabs, array, multirow, xcolor, hyperref, geometry,
%           amsmath, enumitem, tcolorbox, tikz, pgfplots (optional)
% ═══════════════════════════════════════════════════════════════════════

\documentclass[11pt, a4paper]{article}

\usepackage[margin=2.5cm]{geometry}
\usepackage{booktabs}
\usepackage{array}
\usepackage{multirow}
\usepackage{xcolor}
\usepackage{hyperref}
\usepackage{amsmath}
\usepackage{enumitem}
\usepackage{tcolorbox}
\usepackage{parskip}
\usepackage{microtype}
\usepackage{titlesec}
\usepackage{tabularx}

% ── Colour palette ──────────────────────────────────────────────────
\definecolor{opera}{HTML}{2563EB}       % OPERA blue
\definecolor{muted}{HTML}{6B7280}       % grey annotations
\definecolor{positive}{HTML}{059669}    % green — expected advantage
\definecolor{negative}{HTML}{DC2626}    % red — expected disadvantage
\definecolor{neutral}{HTML}{D97706}     % amber — uncertain

\tcbset{
  planbox/.style={
    colback=opera!5, colframe=opera!60,
    fonttitle=\bfseries, left=4pt, right=4pt, top=4pt, bottom=4pt,
  },
  notebox/.style={
    colback=neutral!8, colframe=neutral!50,
    fonttitle=\bfseries\small, left=4pt, right=4pt, top=2pt, bottom=2pt,
  },
}

\hypersetup{colorlinks=true, linkcolor=opera, citecolor=muted, urlcolor=opera}
\titleformat{\section}{\large\bfseries}{}{0em}{}[\vspace{-0.3em}\rule{\textwidth}{0.4pt}]
\titleformat{\subsection}{\normalsize\bfseries\color{opera}}{}{0em}{}

% ── Shorthand ────────────────────────────────────────────────────────
\newcommand{\op}{\textsc{opera}}
\newcommand{\yes}{\textcolor{positive}{\textbf{+}}}
\newcommand{\no}{\textcolor{negative}{\textbf{--}}}
\newcommand{\maybe}{\textcolor{neutral}{\textbf{?}}}
\newcommand{\code}[1]{\texttt{\small #1}}

% ════════════════════════════════════════════════════════════════════
\begin{document}

\begin{center}
  {\LARGE \textbf{OPERA Analysis Plan}} \\[0.4em]
  {\large Outcome-guided Patient Embedding for Registry Adaptation} \\[0.2em]
  {\color{muted}\small Working document — \today}
\end{center}

\vspace{0.5em}

% ════════════════════════════════════════════════════════════════════
\section{Scientific Question}

\begin{tcolorbox}[planbox, title=Core question]
How should a general EHR foundation model be adapted to a specialised
haematology registry to maximise outcome prediction across multiple
diseases and outcomes simultaneously — and does joint multi-disease
training produce a single model that outperforms disease-specific models?
\end{tcolorbox}

\medskip

The question is motivated by the DALY-CARE research programme:

\begin{enumerate}[leftmargin=2em, itemsep=2pt]
  \item \textbf{DALY-CARE} established the infrastructure for observing patients across haematological malignancies in a unified registry.
  \item \textbf{HemaSphere paper} showed empirically that joint training across lymphoma subtypes improves prediction — diseases share structure.
  \item \textbf{OPERA} asks: can we make this cross-disease sharing principled and measurable, using outcomes as supervision signals?
\end{enumerate}

The key intellectual claim is that \emph{outcome trajectories encode
clinical knowledge that clinical presentation alone does not capture},
and that a foundation model adapted with this signal will generalise
better than any tabular model trained on disease-specific features.

% ════════════════════════════════════════════════════════════════════
\section{Experimental Design Overview}

The experiment is a full factorial over three axes:

\medskip
\begin{center}
\begin{tabular}{lll}
  \toprule
  \textbf{Axis} & \textbf{Levels} & \textbf{Script} \\
  \midrule
  Cohort scope  & Per-cohort vs. all cohorts jointly & \code{finetune.py} / \code{joint\_finetune.py} \\
  Model variant & 5 tiers (see \S\ref{sec:variants}) & \code{sweep.py} \\
  Outcome       & 5 outcomes × 2--3 windows & \code{sweep.py} \\
  \bottomrule
\end{tabular}
\end{center}

\medskip
All results flow through \code{opera/run/sweep.py}, which produces a
single \code{results\_table.csv} and \code{results\_table\_auroc.tex}
directly usable in the paper.

% ════════════════════════════════════════════════════════════════════
\section{Cohorts and Outcomes}
\label{sec:cohorts}

\subsection{Disease Cohorts}

\begin{center}
\begin{tabular}{lll}
  \toprule
  \textbf{Cohort} & \textbf{Abbreviation} & \textbf{Notes} \\
  \midrule
  Diffuse large B-cell lymphoma & DLBCL & Primary benchmark (HemaSphere) \\
  Chronic lymphocytic leukaemia & CLL   & Slow natural history \\
  Myelodysplastic syndrome      & MDS   & Competing risks prominent \\
  Multiple myeloma              & MM    & Heavy treatment heterogeneity \\
  Follicular lymphoma           & FL    & Long follow-up \\
  \bottomrule
\end{tabular}
\end{center}

\subsection{Outcomes}

\begin{center}
\begin{tabularx}{\textwidth}{llXl}
  \toprule
  \textbf{Outcome} & \textbf{Windows} & \textbf{Rationale} & \textbf{Cohorts} \\
  \midrule
  All-cause mortality       & 1y, 2y       & Primary survival endpoint; clean signal  & All \\
  Treatment failure         & 6m, 1y, 2y   & Disease control; DLBCL primary target   & DLBCL, CLL, MM \\
  Acute kidney injury       & 7d, 30d      & Treatment toxicity; shared mechanism    & All \\
  Severe infection          & 30d, 90d     & Immunosuppression; cross-disease signal & All \\
  Adverse events (composite)& 30d, 90d     & Safety profile                          & All \\
  \bottomrule
\end{tabularx}
\end{center}

\begin{tcolorbox}[notebox, title=Note on IPI baselines]
  For each cohort where an IPI-family score is available (NCCN-IPI for
  DLBCL, FLIPI-2 for FL, CLL-IPI for CLL), compute AUROC from the
  pre-computed score column in the population CSV.  This is a
  \emph{competing baseline}, not a training input.  Report as a
  separate column in every results table.
\end{tcolorbox}

% ════════════════════════════════════════════════════════════════════
\section{Model Variants}
\label{sec:variants}

The experimental design has four tiers, designed so that each tier
answers a specific question by comparison with the tier below.

\medskip
\begin{tcolorbox}[planbox, title=Tier structure and contrasts of interest]
\begin{description}[leftmargin=0em, itemsep=4pt]
  \item[Tier 0 — Clinical baseline (IPI)]
    NCCN-IPI / CLL-IPI / FLIPI-2 computed from RKKP variables.
    \emph{Contrast}: Does any ML model beat the clinical standard?

  \item[Tier 1 — Tabular ceiling (with RKKP)]
    XGBoost on 87 features including RKKP quality registry variables.
    Trained per cohort per outcome.
    \emph{Contrast}: What is the ceiling when all structured data is available?

  \item[Tier 2 — Tabular floor (EHR only)]
    XGBoost on EHR-derived features only, no RKKP.
    Trained per cohort per outcome.
    \emph{Contrast}: What does a foundation model need to beat to justify its complexity?

  \item[Tier 3 — Foundation model variants]
    Shared encoder, evaluated under four adaptation conditions:
    \begin{enumerate}[leftmargin=1.5em, itemsep=2pt, label=(\alph*)]
      \item \textbf{Pretrain only} — general 3M-patient encoder, no
            haematology-specific adaptation. Lower bound for FMs.
      \item \textbf{DAPT} — continued MLM pretraining on DALY-CARE.
            Tests domain shift correction alone.
      \item \textbf{OPERA (contrastive, per-cohort)} — DAPT + survival-informed
            multi-outcome contrastive adaptation. Tests outcome-guided
            representation learning within a single disease.
      \item \textbf{OPERA (contrastive, joint)} — same model trained on
            all cohorts simultaneously. Tests cross-disease transfer.
    \end{enumerate}

  \item[Tier 4 — Joint multi-task finetuning (bonus)]
    Single model finetuned on all cohorts and all outcomes jointly.
    Evaluated per cohort × outcome cell.
    \emph{Contrast}: Does joint finetuning further improve on separately
    finetuned OPERA? This is where the foundation model's structural
    advantage over tabular models is most direct.
\end{description}
\end{tcolorbox}

\subsection{Why each contrast matters}

\begin{center}
\begin{tabularx}{\textwidth}{lXl}
  \toprule
  \textbf{Comparison} & \textbf{What it answers} & \textbf{Expected direction} \\
  \midrule
  DAPT $>$ Pretrain only
    & Is haematology domain shift real and correctable?
    & \yes \\
  OPERA $>$ DAPT
    & Do outcomes add signal beyond what sequence distributions encode?
    & \yes \\
  OPERA joint $>$ OPERA per-cohort
    & Does cross-disease contrastive training help?
    & \yes\ (especially small cohorts) \\
  Joint finetune $>$ per-cohort finetune
    & Does a single joint model beat separately trained models?
    & \yes\ for FM; \no\ for tabular \\
  FM EHR-only $>$ Tabular EHR-only
    & Is the FM better at using the same data?
    & \yes\ (our primary bet) \\
  Tabular RKKP $>$ FM EHR-only
    & Does structured RKKP data close the gap?
    & \yes\ (expected ceiling effect) \\
  FM + RKKP hybrid $>$ Tabular RKKP
    & Does the FM improve even when RKKP is available?
    & \maybe \\
  IPI $<$ any ML
    & Is ML clinically useful at all?
    & \yes \\
  \bottomrule
\end{tabularx}
\end{center}

% ════════════════════════════════════════════════════════════════════
\section{The Sigma Analysis}
\label{sec:sigma}

The learned $\sigma_k$ parameters from the contrastive stage are a
scientific finding independent of prediction performance.

\begin{tcolorbox}[planbox, title=Sigma interpretation]
  Low $\sigma_k$ $\rightarrow$ outcome $k$ strongly structures the
  embedding space. High $\sigma_k$ $\rightarrow$ noisy / hard to learn jointly.

  \medskip
  Testable prediction: $\sigma_k$ values from the contrastive stage
  should correlate with $\sigma_k$ values from the joint finetune stage
  (both using Kendall weighting). Concordance validates that the
  contrastive stage learned a clinically coherent representation, not
  just an optimisation artefact.
\end{tcolorbox}

The sigma analysis is a publishable result regardless of whether AUROC
improves. If mortality gets low $\sigma$ and adverse events get high
$\sigma$, that is a finding about the latent structure of haematological
disease that no tabular paper can produce.

% ════════════════════════════════════════════════════════════════════
\section{Uncertainty and Clinician Validation}
\label{sec:uncertainty}

\begin{tcolorbox}[planbox, title=Clinician validation protocol]
  \begin{enumerate}[leftmargin=1.5em, itemsep=3pt]
    \item Run \code{extract\_uncertainty()} with $N=30$ MC Dropout samples
          on the held-out test set.
    \item Extract the top-100 most uncertain patients using
          \code{rank\_by\_uncertainty(results, top\_k=100)}.
    \item Present each case to 2--3 haematologists (blinded to model output).
    \item Ask: \emph{``Would you want additional information before making
          a treatment decision?''} (binary, ecologically valid).
    \item Compute point-biserial correlation between model uncertainty
          and clinician uncertainty ratings.
  \end{enumerate}

  \medskip
  \textbf{If correlation is strong:} model uncertainty identifies
  clinically ambiguous cases. This is a usability finding — the model
  knows when to say ``I need more information.''

  \textbf{If correlation is weak:} model is uncertain for statistical
  reasons that do not map to clinical difficulty. Worth reporting honestly.
\end{tcolorbox}

% ════════════════════════════════════════════════════════════════════
\section{Tables for the Paper}

\subsection{Table 1 — Main results (per cohort, primary outcome)}

One row per disease cohort. Primary outcome is all-cause mortality at
1 year. Columns are model variants, metric is AUROC with 95\% bootstrap CI.

\medskip
\begin{tcolorbox}[notebox]
  \textbf{Script:} \code{sweep.py --config sweep\_example.yaml}
  produces \code{results\_table\_auroc.tex} directly. \\
  \textbf{Filter to:} \code{outcome == "mortality\_1y"}, all cohorts, all variants.
\end{tcolorbox}

\medskip
\begin{center}
\begin{tabular}{lccccccc}
  \toprule
  & \multicolumn{2}{c}{\textit{Clinical}} & \multicolumn{2}{c}{\textit{Tabular}} & \multicolumn{3}{c}{\textit{Foundation model}} \\
  \cmidrule(lr){2-3}\cmidrule(lr){4-5}\cmidrule(lr){6-8}
  \textbf{Cohort} & IPI & — & +RKKP & EHR only & Pretrain & DAPT & \op \\
  \midrule
  DLBCL  & & & & & & & \\
  CLL    & & & & & & & \\
  MDS    & & & & & & & \\
  MM     & & & & & & & \\
  FL     & & & & & & & \\
  \midrule
  \textit{Mean} & & & & & & & \\
  \bottomrule
  \multicolumn{8}{l}{\small AUROC (95\% bootstrap CI). Bold = best per row.}
\end{tabular}
\end{center}

\subsection{Table 2 — Multi-outcome results (DLBCL)}

One row per outcome and time window for the primary cohort (DLBCL).
All model variants. Illustrates that the advantage pattern is
outcome-dependent.

\medskip
\begin{tcolorbox}[notebox]
  \textbf{Filter to:} \code{cohort == "dlbcl"}, all outcomes × windows, all variants.
\end{tcolorbox}

\subsection{Table 3 — Joint vs.\ per-cohort (all cohorts, all outcomes)}

Directly tests the collaborator's hypothesis: does a single joint model
outperform separately trained models?

\medskip
Rows: cohort × outcome. Columns: (a) per-cohort OPERA finetune,
(b) joint OPERA finetune, (c) per-cohort tabular RKKP,
(d) joint tabular (XGBoost with disease dummies). Metric: AUROC.

\medskip
\begin{tcolorbox}[notebox]
  Key comparison: column (b) vs.\ column (a) tests the foundation
  model's joint-training hypothesis. Column (d) vs.\ column (c) tests
  the same for tabular — expected to be flat or negative. The
  \emph{interaction} is the finding.
\end{tcolorbox}

\subsection{Table 4 — Sigma values across outcomes}

The learned $\sigma_k$ values from (a) the contrastive stage and
(b) the joint finetune stage.

\begin{center}
\begin{tabular}{lcc}
  \toprule
  \textbf{Outcome} & $\sigma_k$ (contrastive) & $\sigma_k$ (joint finetune) \\
  \midrule
  Mortality 1y        & & \\
  Treatment failure 1y & & \\
  AKI 30d             & & \\
  Severe infection 90d & & \\
  Adverse events 90d  & & \\
  \bottomrule
  \multicolumn{3}{l}{\small Lower = outcome strongly structures the embedding space.}
\end{tabular}
\end{center}

% ════════════════════════════════════════════════════════════════════
\section{Figures for the Paper}

\subsection{Figure 1 — Architecture overview}

A three-panel diagram showing the full OPERA pipeline:
\begin{enumerate}[leftmargin=2em, itemsep=2pt]
  \item \textbf{Left}: General pretraining on $\sim$3M patients (MLM).
  \item \textbf{Centre}: DAPT on DALY-CARE haematology cohort +
        OPERA contrastive stage with survival-informed soft loss and
        DAPT-prior cohort similarity weights.
  \item \textbf{Right}: Per-outcome finetuning and evaluation.
\end{enumerate}
Annotate with: token prefixes, index date concept, outcome windows,
$\sigma_k$ position in the loss.

\medskip
\begin{tcolorbox}[notebox]
  \textbf{Source:} Draw in Excalidraw / Inkscape. Export as PDF vector.
  Use existing architecture sketch from previous sessions as base.
\end{tcolorbox}

\subsection{Figure 2 — Survival pair weighting illustration}

A 2×2 patient grid illustrating how pair weights are computed:
\begin{itemize}[leftmargin=2em, itemsep=1pt]
  \item Axes: time similarity (x) × censoring reliability (y).
  \item Four quadrants: both observed + similar times (weight $\approx$1),
        both observed + distant times (weight low),
        one censored (weight $\approx$0.5), both censored (weight $\approx$0.1).
  \item Overlay: DAPT-prior modulation shifts all weights by cohort similarity.
\end{itemize}

\subsection{Figure 3 — Embedding space shift (UMAP)}

Two-panel UMAP of patient embeddings, coloured by 1-year mortality:
\begin{itemize}[leftmargin=2em, itemsep=1pt]
  \item \textbf{Left}: DAPT-stage embeddings (before contrastive training).
  \item \textbf{Right}: OPERA-stage embeddings (after contrastive training).
\end{itemize}
Show all cohorts in the same plot (use shape for disease, colour for
outcome). The expected finding: OPERA separates outcome-similar patients
across disease boundaries that DAPT does not.

\medskip
\begin{tcolorbox}[notebox]
  \textbf{Script:} \code{opera/visualization/embedding\_plots.py} —
  \code{plot\_embedding\_projection(embeddings, labels, method="umap")}. \\
  Run on held-out test set using DAPT and OPERA checkpoints.
\end{tcolorbox}

\subsection{Figure 4 — Sigma heatmap}

Heatmap: rows = outcomes, columns = cohorts.
Cell value = $\sigma_k$ learned by the multi-cohort OPERA model.

The expected pattern: mortality has low (consistent) $\sigma$ across
cohorts; treatment failure has high variance in $\sigma$ across cohorts
(disease-specific signal). This is the ``what the model learned''
figure.

\subsection{Figure 5 — Joint vs.\ per-cohort performance delta}

A dot plot (Cleveland-style) showing $\Delta$AUROC (joint $-$ per-cohort)
for each cohort × outcome cell.

\begin{itemize}[leftmargin=2em, itemsep=1pt]
  \item Foundation model variants: expected mostly positive deltas.
  \item Tabular variants: expected deltas centred near zero or negative.
  \item Points above the zero line confirm the collaborator's hypothesis.
\end{itemize}

\medskip
\begin{tcolorbox}[notebox]
  The interaction here is the paper's clearest demonstration of why
  foundation models are structurally better suited to multi-disease
  registries. Use a side-by-side panel: FM delta (left) vs.\ tabular
  delta (right).
\end{tcolorbox}

\subsection{Figure 6 — Clinician uncertainty validation}

Scatter plot: model uncertainty (x-axis, embedding variance from MC
Dropout) vs.\ clinician uncertainty rating (y-axis, proportion of
reviewers requesting additional information).

Annotate with: Spearman $\rho$, $p$-value, $n$ cases.

Expected shape: positive correlation. Highlight 3--5 representative
cases in the margin with brief clinical descriptions.

\subsection{Figure 7 — Label efficiency curves (supplementary)}

Learning curves showing AUROC vs.\ number of labelled training
examples for (a) OPERA encoder and (b) tabular XGBoost, evaluated on
DLBCL mortality.

The expected finding: OPERA reaches tabular performance with $\sim$10--20\%
of the labelled data. This quantifies the label efficiency advantage of
pretraining and is a strong practical argument for the approach.

\medskip
\begin{tcolorbox}[notebox]
  \textbf{How to generate:} subsample training set to $[5\%, 10\%, 20\%,
  40\%, 60\%, 80\%, 100\%]$ of labels. Run \code{finetune.py} for each
  fraction. Plot AUROC vs.\ fraction, with 95\% CI from 3--5 seeds.
\end{tcolorbox}

% ════════════════════════════════════════════════════════════════════
\section{Analysis Checklist}

\begin{center}
\begin{tabularx}{\textwidth}{clXl}
  \toprule
  & \textbf{Step} & \textbf{Description} & \textbf{Status} \\
  \midrule
  \square & DAPT embedding store
    & Run \code{build\_dapt\_embedding\_store()} once after DAPT checkpoint
    & $\square$ \\
  \square & Outcome parquets
    & Add \code{time\_days} and \code{event} columns to all outcome parquets
    & $\square$ \\
  \square & IPI scores
    & Compute IPI per cohort, add to \code{population\_full.csv}
    & $\square$ \\
  \square & Per-cohort contrastive
    & Run \code{contrastive.py} per cohort (DLBCL first)
    & $\square$ \\
  \square & Multi-cohort contrastive
    & Run \code{contrastive\_multicohort.py} on all cohorts jointly
    & $\square$ \\
  \square & Per-cohort finetune sweep
    & Run \code{sweep.py} with per-cohort OPERA variant
    & $\square$ \\
  \square & Joint finetune
    & Run \code{joint\_finetune.py} — single outcome first, then multi-output
    & $\square$ \\
  \square & Joint finetune sweep
    & Add \code{opera\_joint} variant to sweep, rerun
    & $\square$ \\
  \square & Tabular baselines
    & Train XGBoost per cohort × outcome, export \code{results/tabular/*.json}
    & $\square$ \\
  \square & Sigma analysis
    & Extract and compare $\sigma_k$ from contrastive and joint finetune
    & $\square$ \\
  \square & Uncertainty extraction
    & Run \code{extract\_uncertainty()} on held-out set, $N=30$ passes
    & $\square$ \\
  \square & Clinician validation
    & Present top-100 uncertain cases to haematologists
    & $\square$ \\
  \square & UMAP plots
    & Generate Fig.\ 3 from DAPT and OPERA checkpoints
    & $\square$ \\
  \square & Label efficiency
    & Run subsampled finetunes for Fig.\ 7
    & $\square$ \\
  \bottomrule
\end{tabularx}
\end{center}

% ════════════════════════════════════════════════════════════════════
\section{Expected Results and Narrative}

The paper's argument structure, in order of evidence:

\begin{enumerate}[leftmargin=2em, itemsep=6pt]

  \item \textbf{FM beats IPI} across all cohorts and primary outcomes.
        Establishes clinical relevance. (Table 1)

  \item \textbf{OPERA $>$ DAPT $>$ Pretrain only} within each cohort.
        Establishes that each adaptation stage adds something.
        The DAPT $\to$ OPERA gap is our main methodological contribution.
        (Table 1, all cohorts)

  \item \textbf{FM EHR-only approaches tabular RKKP}.
        The gap should be small or close on several outcomes.
        This is the strong version of the claim: the FM reconstructs
        RKKP-equivalent structure from longitudinal EHR alone.
        (Replication of HemaSphere result in FM setting)

  \item \textbf{Joint training helps FMs, not tabular models}.
        The $\Delta$AUROC (joint $-$ per-cohort) is positive for FM
        variants and near-zero or negative for XGBoost.
        This is the collaborator's hypothesis and the paper's novel
        empirical finding. (Figure 5, Table 3)

  \item \textbf{$\sigma_k$ values are clinically coherent}.
        Mortality and AKI get low $\sigma$ (clean, shared signal);
        treatment failure gets higher $\sigma$ (disease-specific).
        Concordance across contrastive and joint finetune stages
        validates the representation. (Table 4, Figure 4)

  \item \textbf{Uncertainty tracks clinical difficulty}.
        Positive correlation between model uncertainty and clinician
        uncertainty ratings validates that the model has learned what
        makes a case hard. (Figure 6)

\end{enumerate}

\begin{tcolorbox}[planbox, title=If results are strong enough for Nature Machine Intelligence]
  The framing shifts from ``we adapted a foundation model'' to:

  \medskip
  \textit{``Foundation models pretrained on general EHR data are
  uniquely suited to joint multi-disease clinical registries because
  their shared token space naturally encodes cross-disease similarity
  structure that tabular models require explicit feature engineering to
  approximate — and outcome-guided adaptation makes this structure
  measurable and scientifically interpretable.''}

  \medskip
  This requires: (a) the joint FM advantage over tabular is consistent
  and substantial; (b) the $\sigma_k$ analysis is interpretable; (c) the
  clinician uncertainty result holds. Any two of three is sufficient.
\end{tcolorbox}

% ════════════════════════════════════════════════════════════════════
\section{Anticipated Reviewer Challenges}

\begin{center}
\begin{tabularx}{\textwidth}{lX}
  \toprule
  \textbf{Challenge} & \textbf{Response} \\
  \midrule
  ``Why not train from scratch?''
    & Label efficiency curves (Fig.~7) show FM reaches tabular
      performance with $\sim$10--20\% of labelled data. \\[4pt]
  ``DALY-CARE is one registry — does it generalise?''
    & Cross-cohort transfer within registry tests internal
      generalisation. Cross-disease sigma concordance validates
      biological coherence. \\[4pt]
  ``Contrastive learning is not novel''
    & Survival-informed soft pair weighting with censoring-aware
      reliability is novel. DAPT-prior cross-disease weighting
      is novel. Combination in a clinical registry is novel. \\[4pt]
  ``The tabular ceiling includes RKKP which FM doesn't have''
    & This is a feature, not a bug — it quantifies exactly what
      the FM needs to close. The hybrid model tests whether
      giving the FM RKKP closes the gap. \\[4pt]
  ``How do you handle censoring?''
    & Reliability weights (1.0 / 0.5 / 0.1) based on whether
      ordering is known. Only pairs where at least one event
      is observed contribute full signal. Report censoring
      fraction per cohort × outcome as a sensitivity table. \\
  \bottomrule
\end{tabularx}
\end{center}

\end{document}
